\documentclass[14pt,a4paper]{extarticle}
\usepackage[top=20mm,bottom=20mm,left=20mm,right=15mm]{geometry}
\usepackage{fontspec}
\setmainfont{Times New Roman}
\setmonofont{Consolas}[Scale=MatchLowercase]
\usepackage{polyglossia}
\setdefaultlanguage{russian}
\setotherlanguage{english}
\usepackage{booktabs,array,graphicx,titlesec,enumitem}
\usepackage[hyphens]{url}
\usepackage[hidelinks]{hyperref}
\setlength{\parindent}{1.25cm}
\setlength{\parskip}{0pt}
\setlength{\emergencystretch}{2em}
\clubpenalty=10000
\widowpenalty=10000
\sloppy
\urlstyle{same}
\makeatletter
\renewcommand\normalsize{\@setfontsize\normalsize\@xivpt{16.1}}
\makeatother
\AtBeginDocument{\normalsize}
\titleformat{\section}{\centering\bfseries}{}{0pt}{}
\titlespacing*{\section}{0pt}{8pt}{3pt}
\titleformat{\subsection}{\bfseries\itshape}{}{0pt}{}
\titlespacing*{\subsection}{\parindent}{6pt}{2pt}

\hypersetup{pdftitle={Ресурсная эффективность локальных языковых моделей в аграрных задачах: сравнение конфигураций 9B и 27B и использование расчётных инструментов},pdfauthor={Антоний Сарычев; Денис Шестернин}}
\begin{document}
\noindent\textbf{УДК 004.8}\par\medskip
\begin{center}\bfseries\large Ресурсная эффективность локальных языковых моделей в аграрных задачах: сравнение конфигураций 9B и 27B и использование расчётных инструментов\end{center}
\begin{center}\textbf{Антоний Сарычев, Денис Шестернин}\end{center}
\begin{center}\textit{РГАУ-МСХА имени К.А. Тимирязева, Москва, Россия}\end{center}
{\noindent\fontsize{12}{14}\selectfont \textbf{Аннотация.} Локальная модель Qwen 9B дала 141/289 верных ответов на открытом аграрном наборе MERA против 150/289 у 27B и в отдельном замере работала быстрее при меньшей памяти. Выполнение расчётных функций дало обеим моделям 16/16 численных ответов, но ошибки уточнения данных у 9B сохранились.\par}\smallskip
{\noindent\fontsize{12}{14}\selectfont \textbf{Ключевые слова:} языковые модели; локальный вывод; сельское хозяйство; ресурсная эффективность; расчётные инструменты; воспроизводимость.\par}\smallskip
\section*{Введение}
Локальный аграрный ассистент получает запросы двух видов: вопросы на предметные знания и расчёты по явно заданной формуле. Во втором случае связный ответ не гарантирует верного числа, а правильно названный инструмент — того, что расчёт действительно выполнен. При ограниченной памяти персонального компьютера модель приходится выбирать с учётом качества всей конфигурации, задержки ответа, занятой памяти и наблюдавшихся отказов.

Цель работы — описать практический компромисс между качеством и ресурсами у локальных конфигураций Qwen3.5-9B и Qwen3.8-27B и оценить, что даёт реальное выполнение расчётных функций. Рассматриваются три вопроса: (1) какую часть качества 27B сохраняет конфигурация 9B на аграрном наборе и как различаются задержка и память; (2) чем самостоятельный расчёт модели отличается от передачи аргументов проверяемой функции с её фактическим выполнением; (3) дополнительно — что показывают имеющиеся опыты с обращением к облачным моделям, включая отрицательные результаты. Конфигурации различаются поколением, размером, квантованием и размещением в памяти, поэтому причинный эффект числа параметров не выделяется. Конфигурация 9B рассматривается как основной практический вариант, 27B — как вариант для сравнения с более высоким наблюдаемым качеством.

\section*{Связанные работы}
FrugalGPT, AutoMix, Hybrid LLM, RouteLLM и единый подход к маршрутизации и каскадам формализуют выбор модели и передачу запроса более сильной модели [1–5]. RouterBench отделяет анализ сохранённых ответов от поведения работающего маршрутизатора [6]. Опубликованные в этих работах выигрыши не переносятся автоматически на другие модели, данные и оборудование. Исследования семантической энтропии и обзор методов отказа от ответа показывают, что надёжность требует явно заданного критерия неопределённости или отказа [7, 8]; синтаксически корректный JSON или поле самооценки уверенности таким критерием не являются.

PLLaMa, AgriGPT и AgroLLM исследуют адаптацию языковых моделей к аграрной области [9–11], а AgriAgent и AgroTools — планирование и использование инструментов в сельском хозяйстве [12, 13]; сочетание аграрной тематики и инструментов уже сложилось как направление исследований. τ-bench оценивает достигнутый результат взаимодействия агента с пользователем, ToolACE — данные и проверку вызовов функций [14, 15]. Вклад настоящей работы уже: сравнение ресурсов локальных конфигураций и разграничение указанного и фактически выполненного расчёта.

BitNet и технический отчёт Gemma 4 описывают разные подходы к эффективным архитектурам [16, 17], но не содержат измерений для используемых здесь файлов GGUF. AgXQA представляет другой тип аграрных вопросов и использовался только в ранних опытах [18]. Расчётные формулы взяты из руководства FAO-56 [19], платформа оценки MERA описана в [20]. Страница задачи ruTXTAgroBench приведена отдельно как технический источник [Т1].

\section*{Материалы и методы}
\subsection*{Данные}
Набор ruTXTAgroBench платформы MERA содержит 2\,935 русскоязычных вопросов с выбором ответа по десяти аграрным дисциплинам, из них 293 — в открытой части для разработки (DEV) [20, Т1]. Создатели набора сообщают, что вопросы составлены агрономами, отредактированы и перепроверены тремя экспертами; это описание авторов набора, а не наша независимая проверка каждого эталона или пригодности совета для конкретного хозяйства. Вопросы, эталонные ответы и полные журналы хранятся локально с ограничениями на распространение.

Начальная сбалансированная выборка DEV60 входит в DEV293 и использовалась при выборе моделей. Подвыборка TELEMETRY40 (по четыре задания из каждой дисциплины того же DEV, без отбора по правильности) связывает сохранённые ответы с измерением ресурсов; 20 её заданий входят и в DEV60, поэтому независимой проверкой она не является. Опыт C включает 27 сконструированных русскоязычных расчётных заданий с повторяющимися семействами формул. Более ранний опыт HOME150 с синтетическими и двуязычными случаями служит дополнительным контекстом. Точности этих блоков не объединяются; закрытая тестовая часть MERA (TEST) не использовалась.

\subsection*{Оценка качества}
Строгая оценка точного совпадения (EM) требует равенства множеств латинских обозначений выбранных вариантов. F1 — гармоническое среднее точности и полноты множества, усреднённое по заданиям. Четыре исходных эталона смешивают кириллицу и латиницу и не разбираются по заявленному правилу; их не транслитерировали и не исправляли. Основной знаменатель равен 289, все 293 ответа сохранены, консервативный расчёт с нулём для этих четырёх строк приведён отдельно.

Опубликованная процедура оценки MERA использует нормализацию SQuAD, которая, в частности, может удалить латинскую «A» как английский артикль, поэтому её результаты расходятся со строгими EM и F1. Реализацию оценки на сервере MERA независимо проверить не удалось; поэтому воспроизводимая локально строгая метрика приводится отдельно, а официальный результат MERA не заявляется. Интервалы на многократно использованном DEV носят описательный характер; допустимая граница потери качества для 9B заранее не задавалась.

\subsection*{Оборудование и режимы запуска}
Измерения выполнены на одном компьютере: AMD Ryzen 7\,8700F, RTX 5060 Ti с 16\,311 MiB видеопамяти, 31,620 GiB оперативной памяти; Windows 11 Pro, драйвер NVIDIA 617.14, LM Studio 0.4.25+1, среда выполнения llama.cpp CUDA12 версии 2.51.0 (точная сборка движка для каждой модели не раскрывается). Файлы GGUF занимают 8,873 GiB (Qwen9, квантование Q8\_0) и 15,334 GiB (Qwen27, UD-Q4\_K\_M); их идентичность подтверждена контрольными суммами SHA-256.

Исторический прогон DEV293 выполнялся через OpenAI-совместимый интерфейс: строгая JSON-схема ответа, seed 42, контекст 2\,048 токенов, temperature 0, top\_p 1, лимит выходных токенов 512, рассуждения отключены. Запрошенное размещение на GPU — max для Qwen9 и 0,75 для Qwen27, кэширование контрольных точек (checkpoint caching) — 32; фактическое число слоёв на GPU не измерялось.

TELEMETRY40 выполнялась через собственный REST-интерфейс LM Studio (далее — собственный профиль) с теми же входными текстами, контекстом, параметрами декодирования и лимитом токенов, но без явного seed и серверной проверки схемы. Кэширование контрольных точек отключено, размещение на GPU автоматическое, MTP выключен. Порядок блоков ABBA, отдельный синтетический прогрев и явные выгрузки моделей зафиксированы; загрузка и прогрев не входят в задержку запроса.

Память измерялась периодическим опросом в трёх областях: общая занятая видеопамять устройства, общая занятая оперативная память системы и резидентная память (RSS) группы процессов LM Studio; приводятся максимальные наблюдавшиеся значения. Фоновые приложения и разделяемые страницы входят в эти величины, поэтому собственный объём памяти модели не установлен. Задержка — время полного HTTP-запроса на стороне клиента; время до первого токена (TTFT) на стороне клиента и сервера различается. Скорость коротких ответов буквами не характеризует длинную генерацию; энергопотребление не измерялось.

\section*{Результаты}
\subsection*{Качество и ресурсы}
\par\smallskip\noindent\begin{minipage}{\linewidth}\noindent{\fontsize{12}{14}\selectfont Таблица 1. Качество на DEV293 (исторический прогон, OpenAI-совместимый интерфейс). Выполнены все 293 запроса каждой модели; разбираемых эталонов — 289.\par}\nopagebreak\vspace{3pt}
{\centering\fontsize{12}{14}\selectfont \setlength{\tabcolsep}{4pt}
\begin{tabular}{@{}>{\raggedright\arraybackslash}p{0.33\linewidth}>{\centering\arraybackslash}p{0.24\linewidth}>{\centering\arraybackslash}p{0.12\linewidth}>{\centering\arraybackslash}p{0.25\linewidth}@{}}
\toprule
Конфигурация & Строгая EM & F1 & Задержка p50 / p95, с \\
\midrule
Qwen3.5-9B Q8\_0 & 141/289; 48,79 \% & 0,6286 & 0,295 / 0,402 \\
Qwen3.8-27B UD-Q4\_K\_M & 150/289; 51,90 \% & 0,6468 & 1,374 / 1,753 \\
\bottomrule
\end{tabular}\par}\end{minipage}\par\medskip
Разность EM Qwen9 − Qwen27 составляет −3,11 п. п. Обе модели верны в 113 заданиях, только 9B — в 28, только 27B — в 37, обе неверны в 111. Сохранённый бутстреп с расслоением по дисциплинам (10\,000 повторов) даёт описательный 95-процентный интервал разности EM [−8,30; +2,42] п. п. и F1 [−0,06148; +0,02552]. Отношение 141/150 = 94 \% — это отношение числа верных ответов двух конфигураций, а не точность: точность Qwen9 равна 48,79 \%, эквивалентность не установлена. Консервативный расчёт по всем 293 заданиям даёт 141/293 = 48,12 \% и 150/293 = 51,19 \%; истинных ответов четырёх спорных строк он не устанавливает.

\par\smallskip\noindent\begin{minipage}{\linewidth}\noindent{\fontsize{12}{14}\selectfont Таблица 2. TELEMETRY40, собственный профиль: качество сохранённых ответов оценено ретроспективно по 39 эталонам; задержка и память — по 40 завершённым запросам каждой модели.\par}\nopagebreak\vspace{3pt}
{\centering\fontsize{12}{14}\selectfont \setlength{\tabcolsep}{4pt}
\begin{tabular}{@{}>{\raggedright\arraybackslash}p{0.5\linewidth}>{\centering\arraybackslash}p{0.22\linewidth}>{\centering\arraybackslash}p{0.22\linewidth}@{}}
\toprule
Показатель & Qwen9 & Qwen27 \\
\midrule
Строгая EM & 22/39; 56,41 \% & 24/39; 61,54 \% \\
F1 & 0,7197 & 0,7801 \\
Задержка p50 / p95, с & 0,269 / 0,344 & 1,434 / 1,829 \\
TTFT клиента, p50, с & 0,122 & 0,691 \\
Пик видеопамяти устройства, MiB & 11\,471 & 13\,287 \\
Пик оперативной памяти системы, GiB & 20,858 & 28,470 \\
Пик RSS процессов LM Studio, MiB & 11\,198,96 & 18\,742,34 \\
Завершено / нарушения формата JSON & 40/40 / 2 & 40/40 / 0 \\
\bottomrule
\end{tabular}\par}\end{minipage}\par\medskip
На этой подвыборке медианная задержка Qwen9 в 5,338 раза меньше, а Qwen27 верно отвечает на два задания больше; ни одна конфигурация не превосходит другую по всем показателям. Это описательный компромисс без порога допустимой потери, выбранного после наблюдения.

Сверка 80 пар запросов подтвердила побайтовое совпадение системных и пользовательских сообщений и записей об идентичности весов, однако выбранные варианты изменились в 1 из 40 заданий у Qwen9 и в 2 из 40 у Qwen27. Исторические ответы на те же задания дали обеим моделям по 23/39 (F1 0,7282 и 0,7374). Поэтому качество на DEV293 и задержку из TELEMETRY40 нельзя представлять как одновременные измерения при одинаковых настройках; совместно описываются только качество и ресурсы самой TELEMETRY40 — с оговоркой о ретроспективной оценке и повторном использовании DEV.

Первый вариант собственного профиля с кэшированием 32 контрольных точек привёл к остановкам Qwen27 защитным ограничением по оперативной памяти; с отключённым кэшированием TELEMETRY40 завершена полностью, а отказы сохранены в журналах как операционные, а не как ошибки ответа. Одна успешная серия не подтверждает долговременную устойчивость.

\subsection*{Расчётные задачи и выполнение инструментов}
Во всех режимах опыта C модель получает тип задачи, вопрос, входные числа и единицы. Проверяются: сумма эффективных температур (GDD) без верхнего ограничения — max(0, (Tmin + Tmax)/2 − base), суммируемая по заданным дням; давление насыщения и дефицит давления пара (VPD), для одного наблюдения равный давлению насыщения, умноженному на (1 − RH/100), а в суточной версии рассчитываемый по экстремумам температуры и влажности; суточная эталонная эвапотранспирация ET0 по FAO-56 [19]; заданная в опыте разность max(0, target − sum(rain)). Разность осадков не является нормой полива, а ET0 — ни эвапотранспирацией культуры ETc, ни дозой орошения [19]. Недостающие измерения, скорость ветра, поток тепла в почву G и пороги культуры не подставляются ни из эталона, ни из типовых значений.

На историческом этапе опыта C полная правильность Qwen9 и Qwen27 составила 7/27 и 15/27, правильное число с единицами — 1/16 и 6/16, правильно указанный инструмент — 17/26 и 25/26 (неподдерживаемый случай исключён только из последнего знаменателя); расчётный модуль совпал с независимо построенными формульными эталонами в 27/27 случаях. Повторный расчёт по сохранённым ответам дал те же 17/26 и 25/26, но передавал коду исходные данные напрямую и поэтому не проверял передачу аргументов моделью; ответы двух моделей получены при разных профилях кэширования, и их задержки несравнимы.

Новый заранее описанный опыт завершил все 108 решений (27 случаев × 2 модели × 2 режима): 140 измеряемых вызовов моделей и два отдельных синтетических прогрева. Обе модели работали в собственном профиле без кэширования контрольных точек. В режиме без инструментов модель делает один вызов. В режиме с выполнением инструментов модель выбирает разрешённую функцию и передаёт полный набор аргументов; программа сверяет их с исходными данными задания и единицами, выполняет формулу и возвращает результат модели для итогового ответа. Аргументы не дополняются ни из эталона, ни из скрытой категории, ни из ожидаемого инструмента. Это выполнение через JSON на уровне приложения при явно заданных типе задачи и формате обмена; встроенный механизм вызова функций (native function calling) и распознавание намерения в свободном тексте не проверялись.

\par\smallskip\noindent\begin{minipage}{\linewidth}\noindent{\fontsize{12}{14}\selectfont Таблица 3. Опыт C, 27 случаев в каждом режиме. Полная правильность — по исходной процедуре оценки значения, действия и инструмента; соответствие итогового JSON проверено отдельно. Задержка — медиана суммы времени HTTP-запросов к модели без диспетчеризации и журналирования.\par}\nopagebreak\vspace{3pt}
{\centering\fontsize{12}{14}\selectfont \setlength{\tabcolsep}{4pt}
\begin{tabular}{@{}>{\raggedright\arraybackslash}p{0.3\linewidth}>{\centering\arraybackslash}p{0.12\linewidth}>{\centering\arraybackslash}p{0.16\linewidth}>{\centering\arraybackslash}p{0.15\linewidth}>{\centering\arraybackslash}p{0.17\linewidth}@{}}
\toprule
Конфигурация и режим & Верно & Число и единица & Выполнений & Медиана суммы, с \\
\midrule
Qwen9 без инструментов & 7/27 & 1/16 & 0 & 1,469 \\
Qwen9 с инструментами & 18/27 & 16/16 & 16 & 2,857 \\
Qwen27 без инструментов & 15/27 & 6/16 & 0 & 8,989 \\
Qwen27 с инструментами & 25/27 & 16/16 & 16 & 21,879 \\
\bottomrule
\end{tabular}\par}\end{minipage}\par\medskip
Все 32 численных вызова прошли проверку аргументов и дали правильные итоговые значения с единицами. Среди 11 нечисленных случаев Qwen9 получил 6/11 без инструментов и 2/11 с инструментами, Qwen27 — 9/11 в обоих режимах. Следовательно, улучшение арифметики не означает надёжного уточнения недостающих данных и отказа от неподдерживаемой задачи: у Qwen9 четыре таких случая ухудшились. Ошибок связи и усечений по лимиту выходных токенов не наблюдалось. Отдельная проверка полей и типов итогового JSON даёт 24/27 и 26/27 у Qwen9 и 27/27 в обоих режимах у Qwen27, а с требованием разрешённого идентификатора инструмента — 22/27 и 19/27 у Qwen9; баллов в таблице 3 она не меняет, поскольку ответы с нарушением типов уже неверны по исходной оценке.

Лимит 512 выходных токенов задан на один вызов; успешный численный расчёт с инструментом обычно требует двух вызовов и получает вдвое больший суммарный лимит, а обратная связь от инструмента меняется вместе с числом вызовов. Приведена сумма задержек запросов к модели, а не полная задержка приложения. Это один согласованный опыт на материале разработки. Пути и размеры файлов весов сверены с прежней записью; полная контрольная сумма SHA-256 во время замеров повторно не вычислялась.

Сохранены исходные допуски каждого случая (1e-8 для двенадцати высокоточных эталонов; 0,001, 0,005 или 0,01 для четырёх опубликованных округлённых примеров); вспомогательная диагностика с допуском max(исходный допуск, 1e-4) строгих баллов не заменяет. Допуск относится к формуле, а не к агрономической безопасности: совпадение с формульным эталоном не подтверждает правильность полевого решения.

\subsection*{Дополнительные наблюдения: обращение к облачным моделям}
Эти наблюдения не образуют единого рейтинга облачных моделей. Mercury 2.5 оценивалась отдельно на ограниченной выборке MERA DEV60, модели OpenRouter — на собственном наборе HOME DEV9, Codex CLI — как готовый продукт со своей программной обвязкой; результаты Claude Code на заданиях получить не удалось, а модели, размещённые у провайдеров, не обязательно совпадают с локальными файлами GGUF.

Mercury 2.5 (режим instant, 512 токенов) на MERA DEV60 дала 31/59 строгих совпадений, включая четыре неуспешные или усечённые попытки. Повторный расчёт передачи запросов от Qwen9 (37/59) к Mercury по сохранённым ответам не улучшил EM: лучший сохранённый ненулевой порог дал 35/59. Доли обращений к облаку здесь смоделированы, задержка и стоимость каскада не измерялись.

На HOME DEV9 размещённые у провайдера Qwen3.8-27B, Ling3.0-Flash-Sante и Dots3-Note-Preview дали 5/9, 4/9 и 4/9 по процедуре оценки HOME. Codex CLI на MERA DEV60 дал 36/59 при среднем уровне рассуждений и 46 нарушениях формата JSON; задержка этого продукта включает запуск процесса, а подписка имеет стоимость.

В раннем опыте HOME150 режим «инструменты и облако» получил 139/150, локальный каскад — 133/150 при 20 обращениях к облаку вместо 100, но не выполнил заданного критерия качества. В отдельной серии на HOME DEV36 правила дали 31/36, маршрутизатор Jev — 30/36 при росте внешних запросов с 23 до 56. Эти серии показывают, что передача запросов в облако может не окупаться, и подчёркивают необходимость учитывать инструменты и число запросов, но не доказывают пользы облачной маршрутизации для основной задачи.

\section*{Обсуждение и ограничения}
DEV многократно использовался при настройке и выборе моделей; TELEMETRY40 и сконструированный опыт C не являются независимой внешней проверкой. Один компьютер, разные поколения и квантования моделей, разные интерфейсы и размещение в памяти, отключённые рассуждения и короткий лимит выходных токенов ограничивают обобщение. Строгая метрика воспроизводима локально, а реализация оценки на сервере MERA не подтверждена. Экспертное происхождение набора не гарантирует безошибочности эталонов и безопасности рекомендаций.

\section*{Заключение}
В рассмотренном профиле Qwen9 занимает меньше измеренной памяти и отвечает быстрее, а Qwen27 показывает более высокое наблюдаемое качество (150/289 против 141/289 на DEV293). Это основание выбирать конфигурацию под явно заданный практический бюджет, а не вывод об универсальном превосходстве или эквивалентности. Расчётный режим следует оценивать по аргументам, выполнению и итоговому ответу, а не по названию инструмента: реальное выполнение подняло число правильных численных ответов с 1/16 и 6/16 до 16/16 у обеих моделей, но у Qwen9 сохранились ошибки выбора инструмента и уточнения данных, а успешный расчёт требует дополнительного вызова модели. Облачные наблюдения носят дополнительный характер и выигрыша от передачи запросов не показали. Следующий шаг — проверка на независимом тестовом наборе и оценка новых открытых моделей при тех же ограничениях ресурсов.

\section*{Использование ИИ и вклад авторов}
Генеративный ИИ существенно использовался в работе: агенты Codex — при разработке и выполнении экспериментального кода, организации доказательной базы, проверке источников, поддержке анализа и подготовке черновиков; Claude — при редактировании текста и проверке библиографии. Проверка, выполненная ИИ, не является человеческим рецензированием. Авторы несут ответственность за утверждения, библиографию, права и интерпретацию; описание вкладов и окончательное согласие обоих авторов с текстом ожидают подтверждения. Протоколы, контрольные суммы и журналы попыток сохранены в репозитории проекта. Программный код распространяется по лицензии Apache-2.0; лицензия рукописи и права на данные определяются отдельно.

\section*{Список литературы}
\begingroup\sloppy\fontsize{11}{13}\selectfont \begin{enumerate}[label={},leftmargin=0.9cm,itemindent=-0.9cm,itemsep=0pt,parsep=0pt,topsep=2pt]
\item 1. Chen L., Zaharia M., Zou J. FrugalGPT: How to Use Large Language Models While Reducing Cost and Improving Performance // arXiv preprint. 2023. \href{https://arxiv.org/abs/2305.05176}{arXiv:2305.05176}.
\item 2. Aggarwal P., Madaan A., Anand A. et al. AutoMix: Automatically Mixing Language Models // Advances in Neural Information Processing Systems 37 (NeurIPS 2024). 2024. P. 131000–131034. DOI:~\href{https://doi.org/10.52202/079017-4164}{10.52202/079017-4164}.
\item 3. Ding D., Mallick A., Wang C. et al. Hybrid LLM: Cost-Efficient and Quality-Aware Query Routing // International Conference on Learning Representations (ICLR 2024). 2024. \href{https://arxiv.org/abs/2404.14618}{arXiv:2404.14618}.
\item 4. Ong I., Almahairi A., Wu V. et al. RouteLLM: Learning to Route LLMs from Preference Data // International Conference on Learning Representations (ICLR 2025). 2025. Препринт: \href{https://arxiv.org/abs/2406.18665}{arXiv:2406.18665} (под названием «RouteLLM: Learning to Route LLMs with Preference Data»).
\item 5. Dekoninck J., Baader M., Vechev M. A Unified Approach to Routing and Cascading for LLMs // Proceedings of the 42nd International Conference on Machine Learning (ICML 2025). PMLR. 2025. Vol. 267. P. 12987–13010.
\item 6. Hu Q. J., Bieker J., Li X. et al. RouterBench: A Benchmark for Multi-LLM Routing System // arXiv preprint. 2024. \href{https://arxiv.org/abs/2403.12031}{arXiv:2403.12031}.
\item 7. Kossen J., Han J., Razzak M. et al. Semantic Entropy Probes: Robust and Cheap Hallucination Detection in LLMs // arXiv preprint. 2024. \href{https://arxiv.org/abs/2406.15927}{arXiv:2406.15927}.
\item 8. Wen B., Yao J., Feng S. et al. Know Your Limits: A Survey of Abstention in Large Language Models // Transactions of the Association for Computational Linguistics. 2025. Vol. 13. P. 529–556. DOI:~\href{https://doi.org/10.1162/tacl_a_00754}{10.1162/tacl\_a\_00754}.
\item 9. Yang X., Gao J., Xue W. et al. PLLaMa: An Open-source Large Language Model for Plant Science // arXiv preprint. 2024. \href{https://arxiv.org/abs/2401.01600}{arXiv:2401.01600}.
\item 10. Yang B., Zhang Y., Feng L. et al. AgriGPT: a Large Language Model Ecosystem for Agriculture // arXiv preprint. 2025. \href{https://arxiv.org/abs/2508.08632}{arXiv:2508.08632}.
\item 11. Samuel D. J., Skarga-Bandurova I., Sikolia D. et al. AgroLLM: Connecting Farmers and Agricultural Practices through Large Language Models for Enhanced Knowledge Transfer and Practical Application // arXiv preprint. 2025. \href{https://arxiv.org/abs/2503.04788}{arXiv:2503.04788}.
\item 12. Yang B., Zhang Y., Chen Y. et al. AgriAgent: Contract-Driven Planning and Capability-Aware Tool Orchestration in Real-World Agriculture // arXiv preprint. 2026. \href{https://arxiv.org/abs/2601.08308}{arXiv:2601.08308}.
\item 13. Ye Z., Wen Y., Fan X. et al. AgroTools: A Benchmark for Tool-Augmented Multimodal Agents in Agriculture // arXiv preprint. 2026. \href{https://arxiv.org/abs/2605.22366}{arXiv:2605.22366}.
\item 14. Yao S., Shinn N., Razavi P., Narasimhan K. τ-bench: A Benchmark for Tool-Agent-User Interaction in Real-World Domains // International Conference on Learning Representations (ICLR 2025). 2025. \href{https://arxiv.org/abs/2406.12045}{arXiv:2406.12045}.
\item 15. Liu W., Huang X., Zeng X. et al. ToolACE: Winning the Points of LLM Function Calling // International Conference on Learning Representations (ICLR 2025). 2025. \href{https://arxiv.org/abs/2409.00920}{arXiv:2409.00920}.
\item 16. Ma S., Wang H., Ma L. et al. The Era of 1-bit LLMs: All Large Language Models are in 1.58 Bits // arXiv preprint. 2024. \href{https://arxiv.org/abs/2402.17764}{arXiv:2402.17764}.
\item 17. Gemma Team. Gemma 4 Technical Report // arXiv preprint. 2026. \href{https://arxiv.org/abs/2607.02770}{arXiv:2607.02770}.
\item 18. Kpodo J., Kordjamshidi P., Nejadhashemi A. P. AgXQA: A benchmark for advanced Agricultural Extension question answering // Computers and Electronics in Agriculture. 2024. Vol. 225. Art. 109349. DOI:~\href{https://doi.org/10.1016/j.compag.2024.109349}{10.1016/j.compag.2024.109349}.
\item 19. Allen R. G., Pereira L. S., Raes D., Smith M. Crop evapotranspiration: Guidelines for computing crop water requirements. FAO Irrigation and Drainage Paper 56. Rome: FAO, 1998. ISBN 92-5-104219-5.
\item 20. Fenogenova A., Chervyakov A., Martynov N. et al. MERA: A Comprehensive LLM Evaluation in Russian // Proceedings of the 62nd Annual Meeting of the Association for Computational Linguistics (Volume 1: Long Papers). 2024. P. 9920–9948. DOI:~\href{https://doi.org/10.18653/v1/2024.acl-long.534}{10.18653/v1/2024.acl-long.534}.
\end{enumerate}\endgroup
\section*{Технические источники}
\begingroup\sloppy\fontsize{11}{13}\selectfont \begin{enumerate}[label={},leftmargin=0.9cm,itemindent=-0.9cm,itemsep=0pt,parsep=0pt,topsep=2pt]
\item Т1. ruTXTAgroBench: описание задачи MERA Industrial [Электронный ресурс] // MERA. URL: \url{https://mera.a-ai.ru/ru/industrial/tasks/1} (дата обращения: 04.10.2026).
\end{enumerate}\endgroup
\section*{Сведения на английском языке}
{\noindent\fontsize{11}{13}\selectfont \textbf{Title:} Resource-Efficient Local Language Models for Agricultural Tasks: A 9B–27B Comparison with Deterministic Tools\par}\smallskip
{\noindent\fontsize{11}{13}\selectfont \textbf{Authors:} Antoniy Sarychev, Denis Shesternin. Russian State Agrarian University — Moscow Timiryazev Agricultural Academy, Moscow, Russia\par}\smallskip
{\noindent\fontsize{11}{13}\selectfont \textbf{Abstract.} A local Qwen 9B configuration answered 141/289 questions of the public MERA agricultural development set correctly versus 150/289 for 27B and, in a separate measurement, ran faster with less memory. Executing calculation functions gave both models 16/16 numeric answers, but 9B clarification errors remained.\par}\smallskip
{\noindent\fontsize{11}{13}\selectfont \textbf{Keywords:} language models; local inference; agriculture; resource efficiency; calculation tools; reproducibility.\par}\smallskip
\vspace{6pt}{\noindent\fontsize{11}{13}\selectfont \emph{Рабочая редакция от 5 октября 2026 г.; русская и английская версии описывают одно исследование. Согласие авторов, метаданные, лицензия и УДК (рабочий индекс 004.8) подлежат подтверждению.}\par}
\end{document}
